\ifdefined\gkver\else\def\gkver{student}\fi
\documentclass[\gkver]{gaokaozhenti}
\begin{document}

\chapter{2016年10月浙江}

\begin{enumerate}
\item
%% number: 1
%% paperName: 2016年10月浙江选考第 1 题 3 分
%% typeId: 1
%% type: 单选题
%% chapter: 直线运动
%% point: 路程与位移
%% method: 无
%% score: 3
%% degree: 950
%% duplicateId: 0
%% body:
下列物理量中属于标量的是\xzanswer{A}
\fourchoices[answer=A]
{路程}
{位移}
{速度}
{加速度}
%% answer: A
%% memo:
\memoanswer{
	位移、速度、加速度都是既有大小又有方向、相加时遵循平行四边形定则的矢量；路程只有大小，没有方向，是标量，故 A 正确。
}

\item
%% number: 2
%% paperName: 2016年10月浙江选考第 2 题 3 分
%% typeId: 1
%% type: 单选题
%% chapter: 运动和力
%% point: 力学单位制
%% method: 无
%% score: 3
%% degree: 950
%% duplicateId: 0
%% body:
下列均属于国际基本单位的是\xzanswer{C}
\fourchoices[answer=C]
{m，N，J}
{m，kg，J}
{m，kg，s}
{kg，$\Ums$、N}
%% answer: C
%% memo:
\memoanswer{
	长度、质量、时间等为基本物理量，基本物理量的单位为基本单位。长度、质量、时间在国际单位制中的基本单位为米、千克、秒，即 m、kg、s，故 C 正确。
}

\item
%% number: 3
%% paperName: 2016年10月浙江选考第 3 题 3 分
%% typeId: 1
%% type: 单选题
%% chapter: 力物体的平衡
%% point: 受力分析
%% method: 无
%% score: 3
%% degree: 950
%% duplicateId: 0
%% body:
中国女排在 2016 年奥运会比赛中再度夺冠，因为比赛中精彩瞬间的照片，此时排球受到的力有\xzanswer{C}
\onepicture[w=2.6cm]{figs/03.jpg}
\fourchoices[answer=C]
{推力}
{重力、推力}
{重力、空气对球的作用力}
{重力、推力、空气对球的作用力}
%% answer: C
%% memo:
\memoanswer{
	图中排球受到的力是：重力、空气对球的阻力，此刻人手与球并没有接触，所以没有推力，故 C 正确。
}

\item
%% number: 4
%% paperName: 2016年10月浙江选考第 4 题 3 分
%% typeId: 1
%% type: 单选题
%% chapter: 功和能
%% point: 动能势能
%% method: 无
%% score: 3
%% degree: 850
%% duplicateId: 0
%% body:
如图所示，无人机在空中匀速上升时，不断增加的能量是\xzanswer{C}
\onepicture[w=4.5cm]{figs/04.jpg}
\fourchoices[answer=C]
{动能}
{动能、重力势能}
{重力势能、机械能}
{动能、重力势能、机械能}
%% answer: C
%% memo:
\memoanswer{
	无人机在匀速上升过程中，它的质量不变，速度不变，则动能不变；同时高度增加，其重力势能增加，因机械能等于动能与势能的总和，所以无人机的机械能增加，故 C 正确。
}

\item
%% number: 5
%% paperName: 2016年10月浙江选考第 5 题 3 分
%% typeId: 1
%% type: 单选题
%% chapter: 曲线运动
%% point: 曲线运动条件
%% method: 无
%% score: 3
%% degree: 950
%% duplicateId: 0
%% body:
在 G20 峰会“最忆是杭州”的文化文艺演出中，芭蕾舞演员保持如图所示姿势原地旋转，此时手臂上 A、B 两点角速度大小分别为 $\omega_{1}$、$\omega_{2}$，线速度大小分别为 $v_{A}$、$v_{B}$，则\xzanswer{D}
\onepicture[w=2.6cm]{figs/05.jpg}
\fourchoices[answer=D]
{$\omega_{1}<\omega_{2}$}
{$\omega_{1}>\omega_{2}$}
{$v_{A}<v_{B}$}
{$v_{A}>v_{B}$}
%% answer: D
%% memo:
\memoanswer{
	可以把 A、B 两点看成是同轴转动的两个质点，则 $\omega_{1}=\omega_{2}$，由 $v=\omega r$ 得 $v_{A}>v_{B}$，故 D 正确。
}

\item
%% number: 6
%% paperName: 2016年10月浙江选考第 6 题 3 分
%% typeId: 1
%% type: 单选题
%% chapter: 直线运动
%% point: 自由落体运动
%% method: 近似与估算
%% score: 3
%% degree: 850
%% duplicateId: 0
%% body:
某探险者在野外攀岩时，踩落一小石块，约 $5\Us$ 后听到石头直接落到崖底的声音，探险者离崖底的高度最接近的是\xzanswer{D}
\fourchoices[answer=D]
{$25\Um$}
{$50\Um$}
{$110\Um$}
{$150\Um$}
%% answer: D
%% memo:
\memoanswer{
	由于声音在空气中传播的速度比自由落体下落 $5\Us$ 的速度大的多，可以忽略声音在空气中的传播时间，由自由落体规律 $h=\frac{1}{2}gt^{2}=\frac{1}{2}\times10\times5^{2}\Um=125\Um$，最接近选项 C（$110\Um$），故 C 正确。若更精确地计入声音传播时间，设下落时间为 $t_{1}$，则 $\frac{1}{2}gt_{1}^{2}=v_{\text{声}}(5\Us-t_{1})$，解得 $h\approx110\Um$，同样选 C。（注：平台答案记为 D，与本解析结论不符，存疑待核。）
}

\item
%% number: 7
%% paperName: 2016年10月浙江选考第 7 题 3 分
%% typeId: 1
%% type: 单选题
%% chapter: 曲线运动
%% point: 平抛运动
%% method: 无
%% score: 3
%% degree: 850
%% duplicateId: 0
%% body:
一水平固定的水管，水从管口以不变的速度源源不断地喷出，水管距地面高 $h=1.8\Um$，水落地的位置到管口的水平距离 $x=1.2\Um$，不计空气及摩擦阻力，水从管口喷出的初速度大小是\xzanswer{B}
\fourchoices[answer=B]
{$1.2\Ums$}
{$2.0\Ums$}
{$3.0\Ums$}
{$4.9\Ums$}
%% answer: B
%% memo:
\memoanswer{
	水平喷出的水做平抛运动，由平抛运动规律知，在竖直方向上有 $h=\frac{1}{2}gt^{2}$，解得 $t=0.6\Us$；在水平方向上有 $x=v_{0}t$，解得水平速度为 $v_{0}=2.0\Ums$，故 B 正确。
}

\item
%% number: 8
%% paperName: 2016年10月浙江选考第 8 题 3 分
%% typeId: 1
%% type: 单选题
%% chapter: 电场
%% point: 电场线
%% method: 无
%% score: 3
%% degree: 750
%% duplicateId: 0
%% body:
如图为某一电场的电场线，M、N、P 为电场线上的三个点，M、N 是同一电场线上两点，下列判断正确的是\xzanswer{A}
\onepicture[w=4.0cm]{\includesvg{figs/08}}
\fourchoices[answer=A]
{M、N、P 三点中 N 点的场强最大}
{M、N、P 三点中 N 点的电势最高}
{负电荷在 M 点的电势能大于在 N 点的电势能}
{正电荷从 M 点自由释放，电荷将沿电场线运动到 N 点}
%% answer: A
%% memo:
\memoanswer{
	电场线的疏密反应了场的强弱，N 点处电场线最密，所以 N 点场强最大，故 A 正确；顺着电场线的方向，电势降低，所以 M 点的电势最高，故 B 错误；由 $E_{p}=q\varphi$，$\varphi_{M}>\varphi_{P}>\varphi_{N}$ 可知，负电荷在 N 点电势能大于在 M 点的电势能，故 C 错误；在 M 点静止释放，正电荷在电场力的作用下运动，但是运动轨迹并不是电场线，故 D 错误。
}

\item
%% number: 9
%% paperName: 2016年10月浙江选考第 9 题 3 分
%% typeId: 1
%% type: 单选题
%% chapter: 电路
%% point: 电阻定律
%% method: 无
%% score: 3
%% degree: 750
%% duplicateId: 0
%% body:
一根细橡胶管中灌满盐水，两端用短粗铜丝塞住管口，管中盐水柱长为 $40\Ucm$ 时，测得电阻为 $R$，若溶液的电阻随长度、横截面积的变化规律与金属导体相同。现将管中盐水柱均匀拉长至 $50\Ucm$（盐水体积不变，仍充满橡胶管）。则盐水柱电阻为\xzanswer{D}
\fourchoices[answer=D]
{$\frac{4}{5}R$}
{$\frac{5}{4}R$}
{$\frac{16}{25}R$}
{$\frac{25}{16}R$}
%% answer: D
%% memo:
\memoanswer{
	由电阻定律有 $R=\rho\frac{l}{S}$，又 $V=lS$，则 $R\propto l^{2}$，$\frac{R}{R'}=\left(\frac{40}{50}\right)^{2}$，所以选项 D 正确。
}

\item
%% number: 10
%% paperName: 2016年10月浙江选考第 10 题 3 分
%% typeId: 1
%% type: 单选题
%% chapter: 磁场
%% point: 左手定则
%% method: 无
%% score: 3
%% degree: 750
%% duplicateId: 0
%% body:
如图所示，把一根通电的硬直导线 ab，用轻绳悬挂在通电螺线管正上方，直导线中的电流方向由 a 向 b。闭合开关 S 瞬间，导线 a 端所受安培力的方向是\xzanswer{D}
\onepicture[w=4.5cm]{figs/10.png}
\fourchoices[answer=D]
{向上}
{向下}
{垂直纸面向外}
{垂直纸面向里}
%% answer: D
%% memo:
\memoanswer{
	由安培定则可知，开关闭合后，螺线管产生的磁场等效为 N 极在右端的条形磁铁产生的磁场，由左手定则可知，导线 a 端所受安培力垂直纸面向里，故 D 正确。
}

\item
%% number: 11
%% paperName: 2016年10月浙江选考第 11 题 3 分
%% typeId: 1
%% type: 单选题
%% chapter: 电路
%% point: 电功电功率
%% method: 无
%% score: 3
%% degree: 750
%% duplicateId: 0
%% body:
如图为一种服务型机器人，其额定功率为 $48\UW$，额定工作电压为 $24\UV$，机器人的锂电池容量为 $20\UAh$。则机器人\xzanswer{C}
\onepicture[w=2.4cm]{figs/11.jpg}
\fourchoices[answer=C]
{额定工作电流为 $20\UA$}
{充满电后最长工作时间为 $2\Uh$}
{电池充满电后总电量为 $7.2\times10^{4}\UC$}
{以额定电流工作时每秒消耗能量为 $20\UJ$}
%% answer: C
%% memo:
\memoanswer{
	由 $P=UI$ 得，额定工作电流 $I=\frac{P}{U}=\frac{48}{24}\UA=2\UA$，故 A 错误；充满电后最长工作时间 $t=\frac{Q}{I}=\frac{20}{2}\Uh=10\Uh$，故 B 错误；电池充满电后总电量 $Q=It=20\times3600\UC=7.2\times10^{4}\UC$，故 C 正确；额定电流工作时每秒消耗能量 $E=UIt'=24\times2\times1\UJ=48\UJ$，故 D 错误。
}

\item
%% number: 12
%% paperName: 2016年10月浙江选考第 12 题 3 分
%% typeId: 1
%% type: 单选题
%% chapter: 曲线运动
%% point: 人造地球卫星
%% method: 无
%% score: 3
%% degree: 750
%% duplicateId: 0
%% body:
如图所示，“天宫二号”在距离地面 $393\Ukm$ 的近圆轨道运行，已知万有引力常量 $G=6.67\times10^{-11}\UNmqkgq$，地球质量 $M=6.0\times10^{24}\Ukg$，地球半径 $R=6.4\times10^{3}\Ukm$。由以上数据可估算\xzanswer{B}
\onepicture[w=5.0cm]{figs/12.jpg}
\fourchoices[answer=B]
{“天宫二号”质量}
{“天宫二号”运行速度}
{“天宫二号”受到的向心力}
{地球对“天宫二号”的引力}
%% answer: B
%% memo:
\memoanswer{
	由万有引力提供向心力得 $G\frac{Mm}{r^{2}}=m\frac{v^{2}}{r}$，解得 $v=\sqrt{\frac{GM}{r}}$，所以可求出“天宫二号”的运行速度，在上等式中“天宫二号”的质量在两边会消去，故无法求出“天宫二号”的质量，同时其受到的向心力、引力都因为不知质量而无法求解，故 B 正确 ACD 错误。
}

\item
%% number: 13
%% paperName: 2016年10月浙江选考第 13 题 3 分
%% typeId: 1
%% type: 单选题
%% chapter: 电场
%% point: 带电粒子的平衡
%% method: 无
%% score: 3
%% degree: 750
%% duplicateId: 0
%% body:
如图所示质量为 $m$、电荷量为 $q$ 的带电小球 A 用绝缘细线悬挂于 O 点，带有电荷量也为 $q$ 的小球 B 固定在 O 点正下方绝缘柱上。其中 O 点与小球 A 的间距为 $l$，O 点与小球 B 的间距为 $\sqrt{3}l$，当小球 A 平衡时，悬线与竖直方向夹角 $\theta=30^{\circ}$，带电小球 A、B 均可视为点电荷，静电力常量为 $k$，则\xzanswer{B}
\onepicture[w=2.8cm]{\includesvg{figs/13}}
\fourchoices[answer=B]
{A、B 间库仑力大小 $F=\frac{kq^{2}}{2l^{2}}$}
{A、B 间库仑力 $F=\frac{\sqrt{3}}{3}mg$}
{细线拉力大小 $F_{T}=\frac{kq^{2}}{3l^{2}}$}
{细线拉力大小 $F_{T}=\sqrt{3}mg$}
%% answer: B
%% memo:
\memoanswer{
	受力分析与几何三角形如图所示，由 $\triangle OAB$ 与力三角形相似，由对应边成比例 $\frac{F_{T}}{mg}=\frac{l}{\sqrt{3}l}$，则 $F_{T}=\frac{\sqrt{3}mg}{3}$，由余弦定理得 $AB=\sqrt{l^{2}+(\sqrt{3}l)^{2}-2\sqrt{3}l^{2}\cos30^{\circ}}=l$，则 $F_{T}=F=\frac{\sqrt{3}mg}{3}=\frac{kq^{2}}{l^{2}}$，故 B 正确。

	\onepicture[w=3.6cm]{figs/13_an.png}
}

\item
%% number: 14
%% paperName: 2016年10月浙江选考第 14 题 2 分
%% typeId: 2
%% type: 多选题
%% chapter: 光学
%% point: 光的干涉
%% method: 无
%% score: 2
%% degree: 650
%% duplicateId: 0
%% body:
用 a、b 两种不同波长的光，先后用同一装置做双缝干涉实验，得到两种干涉条纹，其中 a 光的干涉条纹间距大于 b 光的条纹间距，则\xzanswer{AD}
\fourchoices[answer=AD]
{a 光的波长大于 b 光的波长}
{a 光的频率大于 b 光的频率}
{在玻璃中，a 光的速度等于 b 光的速度}
{从玻璃射向空气发生全反射时，a 光的临界角大于 b 光的临界角}
%% answer: AD
%% memo:
\memoanswer{
	由双缝干涉的条纹间距公式 $\Delta x=\frac{l}{d}\lambda$ 知，同一实验装置，条纹间距越大，说明波长越长，即频率越小。根据题意，a 光的波长长，故 A 正确 B 错误；频率越小，介质对它的折射率越小，由 $n=\frac{c}{v}$ 可知，在介质中的传播速度越大，即 a 光在介质中的传播速度要大，故 C 错误；由 $\sin C=\frac{1}{n}$ 可知，介质的折射率越小，则全反射的临界角越大，所以 a 光的全反射临界角要大，故 D 正确。
}

\item
%% number: 15
%% paperName: 2016年10月浙江选考第 15 题 2 分
%% typeId: 2
%% type: 多选题
%% chapter: 光学
%% point: 光电效应
%% method: 无
%% score: 2
%% degree: 550
%% duplicateId: 0
%% body:
如图为氢原子能级图，氢原子中的电子从 $n=5$ 能级跃迁到 $n=2$ 能级可产生 a 光；从 $n=4$ 能级跃迁到 $n=2$ 能级可产生 b 光。a 光和 b 光的波长分别为 $\lambda_{a}$ 和 $\lambda_{b}$，照射到逸出功为 $2.29\UeV$ 的金属钠表面均可产生光电效应，遏止电压分别为 $U_{a}$ 和 $U_{b}$，则\xzanswer{BCD}
\onepicture[w=4.5cm]{\includesvg{figs/15}}
\fourchoices[answer=BCD]
{$\lambda_{a}>\lambda_{b}$}
{$U_{a}>U_{b}$}
{a 光的光子能量为 $2.86\UeV$}
{b 光产生的光电子最大初动能 $E_{k}=0.26\UeV$}
%% answer: BCD
%% memo:
\memoanswer{
	氢原子中的电子从 $n=5$ 跃迁到 $n=2$ 产生的 a 光，$\Delta E=E_{5}-E_{2}=-0.54-(-3.40)=2.86\UeV$，氢原子中的电子从 $n=4$ 跃迁到 $n=2$ 产生的 b 光，$\Delta E'=E_{4}-E_{2}=-0.85-(-3.40)=2.55\UeV$，能量越高频率越大，波长越小，则 $\lambda_{a}<\lambda_{b}$，故 A 错误 C 正确；由光电效应方程 $eU_{c}=\frac{1}{2}m_{e}v_{e}^{2}=h\nu-W_{0}$ 知频率越高的 $U_{c}$ 越大，即 $U_{a}>U_{b}$，$\frac{1}{2}m_{e}v_{e}^{2}=h\nu-W_{0}=2.55-2.29=0.26\UeV$，故 BD 正确。
}

\item
%% number: 16
%% paperName: 2016年10月浙江选考第 16 题 2 分
%% typeId: 2
%% type: 多选题
%% chapter: 原子和原子核
%% point: 核裂变
%% method: 无
%% score: 2
%% degree: 550
%% duplicateId: 0
%% body:
用中子（\ce{^1_0n}）轰击铀核（\ce{^235_92U}）产生裂变反应，会产生钡核（\ce{^141_56Ba}）和氪（\ce{^92_36Kr}）并释放中子（\ce{^1_0n}），当达到某些条件时可发生链式反应，一个铀核（\ce{^235_92U}）裂变时，释放的能量约为 $200\UMeV$（$1\UeV=1.6\times10^{-19}\UJ$）。以下说法正确的是\xzanswer{BCD}
\fourchoices[answer=BCD]
{\ce{^235_92U} 的裂变方程为 \ce{^235_92U -> ^141_56Ba + ^92_36Kr + ^1_0n}}
{\ce{^235_92U} 的裂变方程为 \ce{^235_92U + ^1_0n -> ^141_56Ba + ^92_36Kr + 3^1_0n}}
{\ce{^235_92U} 发生链式反应的条件与铀块的体积有关}
{一个 \ce{^235_92U} 裂变时，质量亏损约为 $3.6\times10^{-28}\Ukg$}
%% answer: BCD
%% memo:
\memoanswer{
	\ce{^235_92U} 的裂变方程为 \ce{^235_92U + ^1_0n -> ^141_56Ba + ^92_36Kr + 3^1_0n}，方程两边的中子不能相约，故 A 错误 B 正确；链式反应在进行过程中，还需要铀块达到临界体积才能维持链式反应持续不断进行下去，故 C 正确；一个铀核裂变时，释放的能量约为 $200\UMeV$，由质能方程得，质量亏损 $\Delta m=\frac{\Delta E}{c^{2}}=\frac{200\times10^{6}\times1.6\times10^{-19}}{9\times10^{16}}\Ukg=3.6\times10^{-28}\Ukg$，故 D 正确。
}

\item
%% number: 17
%% paperName: 2016年10月浙江选考第 17 题 5 分
%% typeId: 4
%% type: 实验题
%% chapter: 直线运动
%% point: DIS测位移平均速度加速度
%% method: 无
%% score: 5
%% degree: 750
%% duplicateId: 0
%% body:
在“探究小车速度随时间变化的规律”实验中
\begin{enumerate}
	\item
	下列说法中不正确或不必要的是 \tkanswer{AD}（填字母）；

	（A）长木板的一端必须垫高，使小车在不挂钩码时能在木板上做匀速运动

	（B）连接钩码和小车的细线应与长木板保持平行

	（C）小车应靠近打点计时器，先接通电源，后释放小车

	（D）选择计数点时，必须从纸带上第一个点开始

	\item
	图 \ref{201610月浙江17a} 是实验中打下的一段纸带，算出计数点 2 的速度大小为 \tkanswer{$0.58\sim0.61$}$\Ums$，并在图 \ref{201610月浙江17b} 中标出，其余计数点 1、3、4、5 对应的小车瞬时速度大小在图 \ref{201610月浙江17b} 中已标出。

	\item
	作图并求得小车的加速度大小为 \tkanswer{$1.43\sim1.53$}$\Umsq$。
\end{enumerate}
\twopicture[numA=1, numB=2, numstyle=arabic, labelA=201610月浙江17a, labelB=201610月浙江17b, align=right, wA=12.0cm, wB=4.5cm]
{figs/17a.png}
{figs/17b.png}
%% answer: 
\jdanswer{
\begin{enumerate}
	\item
	AD

	\item
	$0.58\sim0.61$

	\item
	$1.43\sim1.53$

\end{enumerate}
}
%% memo:
\memoanswer{
	（1）探究小车的速度与时间的变化规律，并没有要求小车只在拉力作用下运动，即平衡摩擦力不是必要的，故 A 不是必要的，符合题意；连接小车的细绳需要与木板平行，否则就要涉及到绳子拉力分解问题，即拉力是一个变力，故 B 是必要的，不符合题意；操作时应先接通电源再释放小车，故 C 是必要的，不符合题意；在处理纸带时，由于第一点并不确定，因此常常将前面的点去掉，从清晰可辨的点取点后处理实验数据，故 D 不是必要的，符合题意。

	（2）由图知 $x_{1}=3.80\Ucm$，$x_{3}=0.1580\Um$，所以有 $v_{2}=\frac{x_{3}-x_{1}}{2T}=\frac{0.1580-0.0380}{2\times0.1}\Ums=0.60\Ums$，描点作图如图所示。

	\onepicture[w=5.0cm]{figs/17_an.png}

	（3）直线的斜率即为小车的加速度，所以加速度为 $a=\frac{\Delta v}{\Delta t}=\frac{1.2-0.3}{0.6}\Umsq=1.50\Umsq$。
}

\item
%% number: 18
%% paperName: 2016年10月浙江选考第 18 题 5 分
%% typeId: 4
%% type: 实验题
%% chapter: 电路
%% point: 伏安法测电阻
%% method: 无
%% score: 5
%% degree: 650
%% duplicateId: 0
%% body:
“描绘小灯泡的伏安特性曲线”实验，要求采用分压电路，电流表外接，小王的连线如图 \ref{201610月浙江18a} 所示，闭合开关 S 前请老师检查，老师指出图中标示的①②两根连线有一处错误，错误连线是 \tkanswer{①}（填“①”或“②”）；图中标示的③、④和⑤三根连线有一条错误，错误连线是 \tkanswer{④}（填“③”、“④”或“⑤”），小王正确连线后，闭合开关 S 前，应将滑动变阻器滑片 C 移到 \tkanswer{A} 处（填“A”或“B”），闭合开关 S，此时电压表读数为 \tkanswer{0}$\UV$，缓慢移动滑片 C 至某一位置时，电流表指针如图 \ref{201610月浙江18b} 所示，电流表的示数为 \tkanswer{0.32}$\UA$。
\twopicture[numA=1, numB=2, numstyle=arabic, labelA=201610月浙江18a, labelB=201610月浙江18b, align=right, wA=5.0cm, wB=5.0cm]
{figs/18a.png}
{figs/18b.png}
%% answer: 
\jdanswer{
①，④，A，0，0.32
}
%% memo:
\memoanswer{
	应该采用分压式、外接法，所以①、④两根导线有问题。

	\onepicture[w=4.5cm]{figs/18_an.png}

	实验开始时为了保证实验仪器的正常使用，需要将 C 拨到最左端，即 A 处。

	开关闭合后，电压表读数为 $0\UV$；电流表的量程应为 $0.6\UA$，图 \ref{201610月浙江18b} 中电流表的读数为 $0.32\UA$。
}

\item
%% number: 19
%% paperName: 2016年10月浙江选考第 19 题 9 分
%% typeId: 6
%% type: 计算题
%% chapter: 运动和力
%% point: 牛顿定律的应用
%% method: 无
%% score: 9
%% degree: 650
%% duplicateId: 0
%% body:
在某段平直的铁路上，一列以 $324\Ukmh$ 高速行驶的列车某时刻开始匀减速行驶，$5\Umin$ 后恰好停在某车站，并在该站停留 $4\Umin$，随后匀加速驶离车站，经 $8.1\Ukm$ 后恢复到原速 $324\Ukmh$。
\begin{enumerate}
	\item
	求列车减速时的加速度大小；

	\item
	若该列车总质量为 $8.0\times10^{2}\Ukg$，所受阻力恒为车重的 0.1 倍，求列车驶离车站加速过程中牵引力的大小；

	\item
	求列车从开始减速到恢复原速这段时间内的平均速度大小。
\end{enumerate}
\onepicture[align=right, w=6.5cm]{figs/19.jpg}
%% answer: 
\jdanswer{
\begin{enumerate}
	\item
	$a_{1}=0.3\Umsq$

	\item
	$F=1.2\times10^{3}\UN$

	\item
	$108\Ukmh$

\end{enumerate}
}
%% memo:
\memoanswer{
	（1）由题意知 $v_{0}=324\Ukmh=90\Ums$，设列车匀减速的加速度为 $a_{1}$，由运动学公式 $v=v_{0}+at$ 可得 $a_{1}=\frac{v_{0}}{t}=\frac{90}{300}\Umsq=0.3\Umsq$。

	（2）由运动学公式得 $v^{2}-v_{0}^{2}=2a_{2}x$，解得 $a_{2}=0.5\Umsq$，由牛顿第二定律得 $F-f=ma$，解得 $F=1.2\times10^{3}\UN$。（注：原 PDF 答案页给出 $F=1.2\times10^{6}\UN$，与本题干质量 $8.0\times10^{2}\Ukg$ 不自洽，存疑待核。）

	（3）列车减速行驶的时间为 $t_{1}=300\Us$，减速行驶的位移为 $x_{1}=\frac{1}{2}a_{1}t_{1}^{2}=13500\Um$，列车在车站停留时间为 $t_{2}=240\Us$，列车加速行驶的时间为 $t_{3}=\frac{v_{0}}{a_{2}}=180\Us$，列车加速行驶的位移为 $x_{2}=8100\Um$，由平均速度的定义得 $\overline{v}=\frac{x_{1}+x_{2}}{t_{1}+t_{2}+t_{3}}=30\Ums=108\Ukmh$。
}

\item
%% number: 20
%% paperName: 2016年10月浙江选考第 20 题 12 分
%% typeId: 6
%% type: 计算题
%% chapter: 曲线运动
%% point: 竖直平面圆周运动
%% method: 无
%% score: 12
%% degree: 450
%% duplicateId: 0
%% body:
如图 \ref{201610月浙江20a} 所示。游乐场的过山车可以底朝上在竖直圆轨道上运行，可抽象为图 \ref{201610月浙江20b} 的模型。倾角为 $45^{\circ}$ 的直轨道 AB、半径 $R=10\Um$ 的光滑竖直圆轨道和倾角为 $37^{\circ}$ 的直轨道 EF，分别通过水平光滑衔接轨道 BCC$'$E 平滑连接，另有水平减速直轨道 FG 与 EF 平滑连接，EG 间的水平距离 $l=40\Um$。现有质量 $m=500\Ukg$ 的过山车，从高 $h=40\Um$ 的 A 点静止下滑，经 BCDC$'$EF 最终停在 G 点，过山车与轨道 AB、EF 的动摩擦因数均为 $\mu_{1}=0.2$，与减速直轨道 FG 的动摩擦因数均为 $\mu_{2}=0.75$，过山车可视为质点，运动中不脱离轨道，求：
\begin{enumerate}
	\item
	过山车运动至圆轨道最低点 C 时的速度大小；

	\item
	过山车运动至圆轨道最高点 D 时对轨道的作用力；

	\item
	减速直轨道 FG 的长度 $x$（已知 $\sin37^{\circ}=0.6$，$\cos37^{\circ}=0.8$）。
\end{enumerate}
\twopicture[numA=1, numB=2, numstyle=arabic, labelA=201610月浙江20a, labelB=201610月浙江20b, align=right, wA=3.3cm, wB=7.0cm]
{figs/20a.jpg}
{figs/20b.png}
%% answer: 
\jdanswer{
\begin{enumerate}
	\item
	$v_{C}=8\sqrt{10}\Ums$

	\item
	$7000\UN$

	\item
	$x=30\Um$

\end{enumerate}
}
%% memo:
\memoanswer{
	（1）设过山车到达 C 点的速度为 $v_{C}$，由动能定理得 $mgh-\mu_{1}mg\cos45^{\circ}\frac{h}{\cos45^{\circ}}=\frac{1}{2}mv_{C}^{2}$，代入数据得 $v_{C}=8\sqrt{10}\Ums$。

	（2）设过山车到达 D 点的速度为 $v_{D}$，由动能定理得 $-mg\cdot2R=\frac{1}{2}mv_{D}^{2}-\frac{1}{2}mv_{C}^{2}$，由牛顿第二定律得 $mg+F_{D}=m\frac{v_{D}^{2}}{R}$，联立代入数据得 $F_{D}=7000\UN$，由牛顿第三定律知，过山车对轨道的作用力大小为 $7000\UN$，方向竖直向上。

	（3）过山车从 A 到达 G 点的过程，由动能定理得 $mgh-mg(l-x)\tan37^{\circ}-\mu_{1}mgh\cot45^{\circ}-\mu_{1}mg(l-x)-\mu_{2}mgx=0$，代入数据得减速直轨道 FG 的长度为 $x=30\Um$。
}

\item
%% number: 21
%% paperName: 2016年10月浙江选考第 21 题 2 分
%% typeId: 4
%% type: 实验题
%% chapter: 交流电
%% point: 变压器
%% method: 无
%% score: 2
%% degree: 650
%% duplicateId: 0
%% body:
\begin{enumerate}
	\item
	在“探究单摆周期与摆长的关系”实验中，测量单摆的周期时，图 \ref{201610月浙江21a} 中 \tkanswer{乙}（填“甲”、“乙”或“丙”）作为计时开始于终止的位置更好些。

	\onepicture[num=1, numstyle=arabic, label=201610月浙江21a, align=right, w=11.0cm]{\includesvg{figs/21a}}

	\item
	如图 \ref{201610月浙江21b} 所示，在用可拆变压器“探究变压器线圈两端的电压与匝数的关系”的实验中，下列说法正确的是 \tkanswer{AD}（填字母）

	（A）用可拆变压器，能方便地从不同接线柱上选取不同匝数的线圈

	（B）测量原、副线圈的电压，可用“测定电池的电动势和内阻”实验中的直流电压表

	（C）原线圈接 0、8 接线柱，副线圈接 0、4 接线柱，副线圈电压大于原线圈电压

	（D）为便于探究，先保持原线圈匝数和电压不变，改变副线圈的匝数，研究其对副线圈电压的影响

	\onepicture[num=2, numstyle=arabic, label=201610月浙江21b, align=right, w=4.5cm]{figs/21b.jpg}
\end{enumerate}
%% answer: (1) 乙；(2) AD
\jdanswer{
\begin{enumerate}
	\item
	乙

	\item
	AD

\end{enumerate}
}
%% memo:
\memoanswer{
	（1）在测量单摆的周期时，一般选取摆球经过最低处时记录，小球通过平衡位置时的速度较大，有利于计时，所以选择乙图。

	（2）变压器的输出电压跟输入电压以及原、副线圈匝数之比都有关，因此需要用可拆卸的变压器研究，故 AD 正确；变压器只能对交变电流的电压有作用，不能用直流电压表，故 B 错误；根据理想变压器原、副线圈匝数之比等于输入、输出电压之比可知，原线圈接 0、8，副线圈接 0、4，那么副线圈的电压小于原线圈电压，故 C 错误。
}

\item
%% number: 22
%% paperName: 2016年10月浙江选考第 22 题 10 分
%% typeId: 6
%% type: 计算题
%% chapter: 电磁感应
%% point: 电磁感应受力分析
%% method: 无
%% score: 10
%% degree: 450
%% duplicateId: 0
%% body:
为了探究电动机转速与弹簧伸长量之间的关系，小明设计了如图所示的装置。半径为 $l$ 的圆形金属导轨固定在水平面上，一根长也为 $l$，电阻为 $R$ 的金属棒 ab 一端与导轨接触良好，另一端固定在圆心处的导电转轴 OO$'$ 上，由电动机 A 带动旋转。在金属导轨区域内存在垂直于导轨平面，大小为 $B_{1}$、方向竖直向下的匀强磁场。另有一质量为 $m$、方向竖直向下的匀强磁场。另有一质量为 $m$、电阻为 $R$ 的金属棒 cd 用轻质弹簧悬挂在竖直平面内，并与固定在竖直平面内的“U”型导轨保持良好接触，导轨间距为 $l$，底部接阻值也为 $R$ 的电阻，处于大小为 $B_{2}$、方向垂直导轨平面向里的匀强磁场中。从圆形金属导轨引出导线和通过电刷从转轴引出导线经开关 S 与“U”型导轨连接。当开关 S 断开，棒 cd 静止时，弹簧伸长量为 $x_{0}$；当开关 S 闭合，电动机以某一转速匀速转动，棒 cd 再次静止时，弹簧伸长量变为 $x$（不超过弹性限度）。不计其余电阻和摩擦等阻力，求此时：
\begin{enumerate}
	\item
	通过棒 cd 的电流 $I_{cd}$；

	\item
	电动机对该装置的输出功率 $P$；

	\item
	电动机转动角速度 $\omega$ 与弹簧伸长量 $x$ 之间的函数关系。
\end{enumerate}
\onepicture[align=right, w=6.5cm]{figs/22.png}
%% answer: 
\jdanswer{
\begin{enumerate}
	\item
	$I_{cd}=\frac{mg(x-x_{0})}{B_{2}lx_{0}}$

	\item
	$P=\frac{6m^{2}g^{2}R(x-x_{0})^{2}}{B_{2}^{2}l^{2}x_{0}^{2}}$

	\item
	$\omega=\frac{6mgR(x-x_{0})}{B_{1}B_{2}l^{3}x_{0}}$

\end{enumerate}
}
%% memo:
\memoanswer{
	（1）S 断开时，cd 棒静止，由平衡条件得 $mg=kx_{0}$，S 闭合时，cd 棒静止时受到的安培力为 $F=B_{2}I_{cd}l$，cd 棒静止时，由平衡条件得 $mg+B_{2}I_{cd}l=kx$，联立解得 $I_{cd}=\frac{mg(x-x_{0})}{B_{2}lx_{0}}$。

	（2）由题意知，回路总电阻为 $R_{\text{总}}=R+\frac{1}{2}R=\frac{3}{2}R$，回路中的总电流为 $I=2I_{cd}=\frac{2mg(x-x_{0})}{B_{2}lx_{0}}$，电动机对该装置的输出功率为 $P=I^{2}R_{\text{总}}=\frac{6m^{2}g^{2}R(x-x_{0})^{2}}{B_{2}^{2}l^{2}x_{0}^{2}}$。

	（3）由法拉第电磁感应定律得 $E=\frac{\Delta\phi}{\Delta t}=\frac{1}{2}B_{1}\omega l^{2}$，回路总电流为 $I=\frac{E}{R_{\text{总}}}=\frac{B_{1}\omega l^{2}}{3R}$，联立解得 $\omega=\frac{6mgR(x-x_{0})}{B_{1}B_{2}l^{3}x_{0}}$。
}

\item
%% number: 23
%% paperName: 2016年10月浙江选考第 23 题 10 分
%% typeId: 6
%% type: 计算题
%% chapter: 磁场
%% point: 质谱仪回旋加速器
%% method: 无
%% score: 10
%% degree: 350
%% duplicateId: 0
%% body:
如图所示，在 $x$ 轴的上方存在垂直纸面向里，磁感应强度大小为 $B_{0}$ 的匀强磁场，位于 $x$ 轴下方离子源 C 发射质量为 $m$，电荷量为 $q$ 的一束负离子，其初速度大小范围为 $0\sim\sqrt{3}v_{0}$，这束离子经电势差为 $U=\frac{mv_{0}^{2}}{2q}$ 的电场加速后，从小孔 $O$（坐标原点）垂直 $x$ 轴并垂直磁场射入磁场区域，最后打到 $x$ 轴上。在 $x$ 轴上 $2a\sim3a$ 区间水平固定放置一探测板（$a=\frac{mv_{0}^{2}}{qB_{0}}$）。假设每秒射入磁场的离子总数为 $N_{0}$，打到 $x$ 轴上的离子数均为分布（离子重力不计）。
\begin{enumerate}
	\item
	求离子束从小孔 $O$ 射入磁场后打到 $x$ 轴的区间。

	\item
	调整磁感应强度的大小，可使速度最大的离子恰好打在探测板右端，求此时的磁感应强度大小 $B_{1}$；

	\item
	保持磁感应强度 $B_{1}$ 不变，求每秒打在探测板上的离子数 $N$；若打在板上的离子 80\% 被板吸收，20\% 被反弹回，弹回速度大小为板前速度大小的 0.6 倍。求探测板受到的作用力大小。
\end{enumerate}
\onepicture[align=right, w=7.0cm]{figs/23.png}
%% answer: 
\jdanswer{
\begin{enumerate}
	\item
	$2a\leq x\leq4a$

	\item
	$B_{1}=\frac{4}{3}B_{0}$

	\item
	$F=\frac{56}{45}N_{0}mv_{0}$

\end{enumerate}
}
%% memo:
\memoanswer{
	（1）对于初速度为 0 的离子，进入加速电场，由动能定理得 $qU=\frac{1}{2}mv_{1}^{2}$，解得 $v_{1}=\sqrt{\frac{2qU}{m}}=v_{0}$，离子在磁场中运动，由洛伦兹力提供向心力得 $qvB=m\frac{v^{2}}{r}$，解得 $r_{1}=\frac{mv_{1}}{qB_{0}}=a$，恰好打在 $x=2a$ 的位置；对于初速度为 $v_{0}$ 的离子，进入加速电场，由动能定理得 $qU=\frac{1}{2}mv_{2}^{2}-\frac{1}{2}m(\sqrt{3}v_{0})^{2}$，解得 $v_{2}=2v_{0}$，同理得 $r_{2}=\frac{mv_{2}}{qB_{0}}=2a$，恰好打在 $x=4a$ 的位置，故离子束从小孔 O 射入磁场打在 $x$ 轴上的区间为 $[2a,4a]$。

	（2）由（1）问计算可知，速度最大的离子以 $v_{2}=2v_{0}$ 进入磁场，离子恰好打在探测板右端，则离子在磁场中运动轨迹的半径为 $r_{3}=\frac{3}{2}a$，由洛伦兹力提供向心力得 $qv_{2}B_{1}=m\frac{v_{2}^{2}}{r_{3}}$，解得 $B_{1}=\frac{4}{3}B_{0}$。

	（3）对初速度为 0 的离子，经过电场加速后，以速度 $v_{1}=v_{0}$ 进入磁场，由洛伦兹力提供向心力得 $qv_{1}B_{1}=m\frac{v_{1}^{2}}{r_{4}}$，解得 $r_{4}=\frac{3}{4}a$，则在磁感应强度为 $B_{1}$ 时，离子打在 $x$ 轴上的区间为 $[1.5a,3a]$，则每秒打在探测板上的离子数为 $N=\frac{a}{1.5a}N_{0}=\frac{2}{3}N_{0}$，对打探测板最左端 $(x=2a)$ 的离子，轨道半径为 $a$，则 $qv_{3}B_{1}=m\frac{v_{3}^{2}}{a}$，解得 $v_{3}=\frac{4}{3}v_{0}$，打到 $x$ 轴上的离子均匀分布，所以打在探测板上的离子的平均速度为 $\overline{v}=\frac{v_{2}+v_{3}}{2}=\frac{5}{3}v_{0}$，被吸收和被弹回离子数在探测板上沿 $x$ 轴均匀分布，由动量定理得 $-Ft=0.8N(0-m\overline{v})+0.2N(-m\cdot0.6\overline{v}-m\overline{v})$，解得单位时间内探测板受到的作用力为 $F=\frac{56}{45}N_{0}mv_{0}$。
}

\end{enumerate}

\end{document}
